% Ad Familiares 13.24 (ad-familiares-13-24)
% Source: https://raw.githubusercontent.com/PerseusDL/canonical-latinLit/master/data/phi0474/phi056/phi0474.phi056.perseus-lat1.xml
% Fetched by scripts/fetch_latin.py

% Dateline (Perseus): Scr. Romae, ut videtur, a. 708 (46) .

\ciceroLetterOpener{CICERO SERVIO S.}

\ciceroSection{1} Cum antea capiebam ex officio meo voluptatem, quod memineram quam tibi diligenter Lysonem, hospitem et familiarem meum, commendassem, tum vero, postea quam ex litteris eius cognovi tibi eum falso suspectum fuisse, vehementissime laetatus sum me tam diligentem in eo commendando fuisse. ita enim scripsit ad me, sibi meam commendationem maximo adiumento fuisse, quod ad te delatum diceret sese contra dignitatem tuam Romae de te loqui solitum esse.

\ciceroSection{2} de quo etsi pro tua facilitate et humanitate purgatum se tibi scribit esse, tamen primum, ut debeo, tibi maximas gratias ago, cum tantum litterae meae potuerunt, ut iis lectis omnem offensionem suspicionis, quam habueras de Lysone, deponeres, deinde credas mihi adfirmanti velim me hoc non pro Lysone magis quam pro omnibus scribere, hominem esse neminem, qui umquam mentionem tui sine tua summa laude fecerit ; Lyso vero cum mecum prope cotidie esset unaque viveret, non solum quia libenter me audire arbitrabatur, sed quia libentius ipse loquebatur, omnia mihi tua et facta et dicta laudabat.

\ciceroSection{3} quapropter etsi a te ita tractatur, ut iam non desideret commendationem meam unisque se litteris meis omnia consecutum putet, tamen a te peto in maiorem modum ut eum etiam atque etiam tuis officiis, liberalitate conplectare. scriberem ad te qualis vir esset, ut superioribus litteris feceram, nisi eum iam per se ipsum tibi satis notum esse arbitrarer.
